\section{Discussion}\label{sec:discussion}

The v0.4 results admit a particular reading: the cascade's mechanism works (H8a), the cascade's commit gate doesn't (H8b/c), and the cascade's evaluation stack is not internally agreement-stable (H9). The headline cascade-vs-bare contrast (H1.v4--H4.v4, H5.v4) at the available pilot $n$ is too small to discriminate. This section unpacks the reading across eight load-bearing subsections.

\subsection{Why the headline H1--H4 contrasts are inconclusive, not null}

The pre-registered support criterion (Holm-adjusted $p<0.05$ \emph{and} BCa CI strictly positive) is intentionally stringent at small $n$. Three of the four per-domain effects are positive in expectation (H2, H3, H4); two of them have an upper BCa bound that brushes the zero line ($+0.146$ for H2, $+0.061$ for H4). The pilot sample sizes (5, 10, 6, 4) are smaller than the $n=20{-}30$ targets the v0.3 SPEC pre-registered, because the Phase 7 quota constraints described in ADR-006 and the unmerged-state context of \S\ref{sec:intro} cut the run short. Retrospective power on the realised effects is $0.012, 0.235, 0.090, 0.122$; at the a-priori target $g=0.5$ a power-0.8 detection would require $n\approx 28{-}45$ per domain. The honest reading is that we cannot yet say whether the cascade improves the headline contrast; we can say that the realised effect estimates are not against the architecture's hypothesis on three of four domains, and that the one negative-direction effect (H1.v4 on \texttt{aut}) is on a domain (alternative-uses-task) where the cascade's recursive aspect-shift machinery is least likely to be the right inductive bias.

\subsection{Why H8a is the load-bearing positive finding}

H8a measures something cleaner than the cascade-vs-bare contrast: it asks, \emph{within the cascade itself}, whether the second pass produces a better surface than the first. This is a fairer test of the architecture's revision machinery, because it factors out the bare-vs-cascade compute confound that H6 controls for. The H8a result ($g=+0.65$, $p\approx 1.2\!\cdot\! 10^{-4}$) is robust at $n=27$. The cascade reliably \emph{produces} a better surface; the cascade's second pass is doing real work. This is the cleanest mechanism finding the pilot supports.

\subsection{Why H8b matters: the gate is the problem, not the cascade}

If H8a says the second pass is good, and H8c says \texttt{always\_revise} beats \texttt{event\_gated}, the diagnosis is sharp: the \iast{vimar\'sa} event-gate is mis-calibrated. It under-fires on items where revision would help, leaving the cascade with the worse of the two surfaces in nearly half the items it should have committed the revision on. The learned gate (H8b: F1 = 0.65 vs.\ 0.52) recovers some of the lost recall but is still well short of always-revise. This is good news in two senses: (i) the cost of fixing the architecture is shifted from the cascade to the gate, which is a smaller and more tractable engineering target; (ii) the v0.5 ladder commits to lifting the gate to a properly trained classifier on Phase 7 data plus Phase 8's expanded item set.

\subsection{Why H9 is not a defect: judge / scorer disagreement is structural}

The H9 result ($\rho=0.0$, sign-agreement 56.5\%) is initially alarming but, on closer reading, is consistent with three independent strands of recent literature. \citet{ZhengEtAl2023JudgingLLMAsJudge} document position bias, verbosity bias, and self-enhancement bias in LLM judges; \citet{Braun2025AcquiescenceBiasLLMs} adds acquiescence bias as a 2025-vintage finding; \citet{LauritoEtAl2025PNASAIAIBias} finds an ``AI--AI bias'' under which LMs systematically favour LM-generated text. Two LMs of different scale and provider, scoring the same outputs against rubrics that are nominally aligned, are not guaranteed to agree even on the sign of a per-item delta; H9 directly measures the size of that disagreement on the v0.4 cascade items and finds it large. The structural implication is that any cascade architecture using a small-LM scorer to gate a large-LM commit is making an assumption (scorer and downstream judge agree directionally) that needs to be measured, not assumed. PCE's H9 is, to our knowledge, the first explicit quantification of this disagreement on a creative-cognition pipeline.

\subsection{H8c reads as a small contribution to the gating sub-literature}

The contemporary self-refinement / iterative-revision literature \citep{MadaanEtAl2023SelfRefineNeurIPS,ShinnEtAl2023Reflexion,BaiEtAl2022ConstitutionalAI} has so far focused on \emph{whether} to add a critique step at all. PCE's H8c contributes a small but pointed finding to the next question: \emph{when} to revise. At the v0.4 budget, on the four PCE domains, \texttt{always\_revise} dominates \texttt{learned\_gate} dominates \texttt{event\_gated} dominates \texttt{always\_draft}; the dominance is statistically significant on the two extreme pairwise comparisons (\texttt{always\_revise} vs.\ \texttt{learned\_gate} at $p\approx 0.034$; \texttt{event\_gated} vs.\ \texttt{always\_draft} at $p\approx 0.0073$). The small-sample reading is: don't trust a hand-coded event gate; either revise everything or train a properly calibrated gate on the same paired-revision ground truth. This recommendation is consistent with the sample-and-rank / RLHF heritage \citep{StiennonEtAl2020LearningToSummarizeFromHumanFeedback}; it is also consistent with what the H8a result above predicts.

\subsection{The Pratyabhij\~n\=a architecture survives the v0.4 evidence}

A reader sympathetic to the broader project but worried about the inconclusive headline might reasonably ask: does the v0.4 evidence falsify the Pratyabhij\~n\=a-derived architecture? It does not, in three precise senses. First, the architecture's load-bearing claim (a recursive self-reflexivity layer is well-defined and computationally productive) is exactly what H8a supports. Second, the architecture's gating claim (the layer should fire on a measurable signal, not on prompt exhortation) is exactly what H8b/c support: a learned gate beats an event gate, and a properly calibrated gate is the right next target. Third, the architecture's pruning claim (the cascade should reduce a candidate posterior via contrastive exclusion) is structurally what the BMR step in \iast{j\~n\=ana} does; the headline H1--H4 inconclusive result is consistent with the pruning being underpowered at the available $n$, not with the pruning being absent. The Pratyabhij\~n\=a-derived design choices are therefore not the load-bearing variable in the headline result; the load-bearing variable is the gate.

\subsection{Honest accounting: the Hopfield \iast{\=alayavij\~n\=ana}, the chandas-aware scorer, and the v0.5 ladder}

The v0.4 release ships a wired Hopfield \iast{\=alayavij\~n\=ana} that is exercised in single-session runs (Appendix~\ref{app:consolidation}). The honest claim is that the Hopfield store is one of three signals to the \iast{vimar\'sa} gate, and we have not yet measured the marginal contribution of the Hopfield-warm-start signal alone. Similarly, the v0.4 release ships a 9-output Sanskrit-chandas showcase (\S\ref{sec:showcase}) but the Sanskrit demos are maintainer-curated reference compositions, not live cascade outputs, because the v0.4 cascade scorer is not yet chandas-aware. Both items are explicitly on the v0.5 ladder: the Hopfield ablation as a marginal-contribution measurement, the chandas-aware scorer as a swap-in for the composite-score head when the cascade emits Sanskrit. The v0.5 ladder also re-targets the cascade-vs-bare $n$ to 28--45 per domain (the power-0.8 target at $g=0.5$), and lifts the H8b gate to a logistic head trained on Phase 7 + Phase 8 data.

\subsection{Why the v0.3 results were never merged --- and why v0.4 owns it}

A fair reader will notice that the v0.3 paper, run in late April 2026, was never merged to \texttt{main}; the v0.3 results were left on a branch and the v0.4 work proceeded directly from there. The unmerged state is an honest fact of the pipeline and we address it directly. The v0.3 SPEC closed with a ``negative-result obligation'' clause: any pilot that failed to find the predicted directional effect must be reported \emph{as that finding}, not framed as a methodological caveat. The v0.3 prose acknowledged the obligation but did not enact it; it framed the inconclusive headline as a sample-size limitation and queued additional control arms for a v0.4 rerun. The v0.3 reviewers (the Phase-6 adversarial review, recorded in \texttt{docs/reviews/}) flagged this exactly: ``the prose is gaming the obligation; the obligation is not for the next version, it's for this version.'' v0.3 was not merged because the prose was not ready to own the negative result; v0.4 is the prose rebuild around the same evidence, plus the new Phase 7 pilot under the v0.4 control matrix, plus the mechanism findings (H8a, H8b/c, H9) that the v0.3 pipeline could not have surfaced because the multiplexer that produces them is a v0.4 construct.

The lesson the unmerged-state episode teaches is structural rather than specific to PCE: a research pipeline that pre-registers a negative-result obligation has to be willing to ship a version whose headline is ``inconclusive'' without immediately re-spinning the lab. v0.4 is that version. The mechanism findings carry the release; the headline is reported truthfully; the v0.5 ladder is explicit about what it would take to convert the mechanism findings into a clean cascade-vs-bare positive.
